% ==============================================================================
% APPENDIX A: REPRODUCIBILITY & ARTIFACT DETAILS
% Chair of Wireless Systems (Fachgebiet Drahtlose Systeme)
% ==============================================================================

\chapter{Artifacts and Reproducibility Checklist}
\label{ch:appendix_reproducibility}

To support open science and facilitate experimental reproduction, this appendix provides full details regarding source code repositories, hardware artifacts, and software environments used in this thesis.

\section{Software Repository and Open-Source Code}
\label{sec:app_repo}

The complete source code of the \ac{UTSR} routing pipeline, edge inference scripts, dataset loaders, and benchmark evaluation harnesses is maintained on GitHub:
\begin{itemize}
    \item \textbf{Repository URL}:\\ \url{https://github.com/chair-wireless-systems/utsr-edge-routing}
    \item \textbf{License}: MIT Open Source License
    \item \textbf{Archival DOI}:\\ \url{https://doi.org/10.5281/zenodo.12345678}
\end{itemize}

\section{Software Environment Specifications}
\label{sec:app_environment}

\Cref{tab:software_deps} lists the exact software packages and versions utilized during development and evaluation.

\begin{table}[ht]
    \centering
    \caption{Key software libraries and versions.}
    \label{tab:software_deps}
    \small
    \begin{tabularx}{\textwidth}{llX}
        \toprule
        \textbf{Package / Tool} & \textbf{Version} & \textbf{Purpose} \\
        \midrule
        Ubuntu Linux & 22.04.4 LTS & Host Operating System \\
        Python & 3.11.8 & Primary Programming Language \\
        PyTorch & 2.2.1+cu121 & Deep Learning Framework \\
        HuggingFace Transformers & 4.38.2 & Model Architectures and NLI \\
        ONNX Runtime & 1.17.1 & Edge Inference Acceleration \\
        AWQ (AutoAWQ) & 0.2.4 & 4-bit Model Quantization \\
        SciPy / NumPy & 1.12.0 / 1.26.4 & Mathematical and Statistical Analysis \\
        \bottomrule
    \end{tabularx}
\end{table}

\section{Hardware Wiring and Pinout Configuration}
\label{sec:app_hardware_wiring}

The ultrasonic sensor transducer is wired to the ESP32 microcontroller via GPIO pins 18 (Trigger) and 19 (Echo), operating at a sampling frequency of \SI{40}{\kilo\hertz}. Sensor data is transmitted to the Jetson edge gateway over an encrypted IEEE 802.15.4 serial link.
